% Revision module: mathematical formalization of the benchmark framework.
% Suggested insertion point in sn-article-final.tex:
% after \subsection{Evaluation Framework}, replacing the current one-paragraph
% description of the evaluation workflow.

\subsection{Formal Benchmark Definition}
\label{sec:benchmark_formalization}

The central object of this benchmark is not a single scalar RMSE, but a family of task-conditioned measurement operators applied to the same candidate force field. Let a configuration be denoted by
\[
    x=(\mathbf R,\mathbf Z,\mathbf h)\in\mathcal X,
\]
where $\mathbf R$ contains atomic coordinates, $\mathbf Z$ atomic identities, and $\mathbf h$ the simulation cell. A reference electronic-structure model defines an energy $E_{\mathrm D}(x)$ and force field $\mathbf F_{\mathrm D}(x)=-\nabla_{\mathbf R}E_{\mathrm D}(x)$. A candidate MLFF $M$ defines $E_M(x)$ and $\mathbf F_M(x)=-\nabla_{\mathbf R}E_M(x)$.

\begin{definition}[Equilibrium and task-conditioned evaluation distributions]
\label{def:distributions}
The conventional equilibrium test distribution is
\[
    \mu_{\mathrm{eq}}(x)
    \propto
    \exp[-\beta E_{\mathrm D}(x)]
\]
restricted in practice to relaxed structures, short AIMD trajectories, or random splits of near-equilibrium DFT configurations. Its standard pointwise error is
\[
    \mathcal E_{\mathrm{pt}}(M;\mu_{\mathrm{eq}})
    =
    \left(\mathbb E_{x\sim\mu_{\mathrm{eq}}}
    \left[
    w_E |E_M(x)-E_{\mathrm D}(x)|^2
    +
    w_F \|\mathbf F_M(x)-\mathbf F_{\mathrm D}(x)\|^2
    \right]\right)^{1/2}.
\]
For a physical process $\tau$, such as dopant migration or interlayer sliding, the benchmark distribution is a process-conditioned distribution $\mu_\tau$ supported on a structured set of configurations generated by a path operator. This distribution may place substantial probability mass on strained, high-energy, low-Boltzmann-weight configurations that are rare under $\mu_{\mathrm{eq}}$ but control kinetics and mechanical stability.
\end{definition}

Definition~\ref{def:distributions} makes explicit why equilibrium RMSE can be accurate and still incomplete: it measures the average pointwise error under $\mu_{\mathrm{eq}}$, whereas barrier prediction, diffusion, and sliding depend on $E_M$ and $\nabla E_M$ over regions sampled by $\mu_\tau$. The two tests therefore answer different mathematical questions.

\subsubsection{Pointwise Probes and Process Probes}
\label{sec:point_path_process}

A pointwise probe evaluates a model on isolated configurations,
\[
    \mathcal P_{\mathrm{pt}}(M;\mathcal S)
    =
    \{E_M(x),\mathbf F_M(x):x\in\mathcal S\},
\]
where $\mathcal S$ is a fixed set of DFT-labeled structures. This includes static energy/force RMSE and RDF-preserving near-equilibrium MD checks. A process probe evaluates a functional of a correlated configuration sequence,
\[
    \Gamma = (x_0,x_1,\ldots,x_N),
\]
or of a trajectory distribution generated by $M$. Examples include a migration barrier, endpoint asymmetry, NEB convergence, diffusion coefficient, thermal conductivity, and sliding periodicity. The benchmark therefore measures both local correctness,
\[
    E_M(x_i)\approx E_{\mathrm D}(x_i), \qquad
    \mathbf F_M(x_i)\approx\mathbf F_{\mathrm D}(x_i),
\]
and functional correctness,
\[
    \Phi_M(\Gamma)\approx \Phi_{\mathrm D}(\Gamma),
\]
where $\Phi$ is a task-level observable rather than a pointwise label.

\subsubsection{Migration-Path Selection}
\label{sec:path_selection_formal}

Let $\mathcal A=\{a_1,\ldots,a_K\}$ be metastable dopant sites obtained from relaxed DFT structures or thermally sampled AIMD motifs. A migration task is an ordered pair
\[
    \tau=(a_s,a_t,c),
\]
where $a_s$ and $a_t$ are initial and final sites and $c$ records the chemical and geometric context, such as in-gap motion, penetration into a quintuple layer, or interlayer sliding. The selected task set
\[
    \mathcal T=\mathcal T_{\mathrm{ID}}\cup\mathcal T_{\mathrm{OOD}}\cup\mathcal T_{\mathrm{coll}}
\]
is designed to include: (i) interpolation-like paths whose local environments overlap with AIMD training data, (ii) extrapolation-like paths that pass through strained or high-energy transition states, and (iii) collective displacements where local environments change weakly but long-range registry or periodicity matters.

For each $\tau$, a path is initialized by interpolation or chemically constrained image construction,
\[
    \Gamma^0_\tau=(x_0,\ldots,x_N),\qquad x_0=a_s,\quad x_N=a_t,
\]
and then optimized by a reference or model-specific NEB operator. A path is retained as a benchmark probe only if its endpoints, chemical identity, cell convention, and image ordering are auditable from the raw inputs and outputs.

\subsubsection{Fixed-Path and Self-Relaxed NEB Operators}
\label{sec:neb_operators_formal}

The apparent discrepancy between a fixed-path barrier in Fig.~\ref{fig:barrier} and a self-consistent MLFF-NEB barrier in Fig.~S3 is resolved by defining two different operators.

\paragraph{Reference DFT path.}
For a task $\tau$, the DFT NEB operator returns a converged reference path
\[
    \Gamma_{\mathrm D,\tau}^{*}
    =
    \operatorname{NEB}_{\mathrm D}(x_0,x_N)
    =
    (x_{0,\tau}^{\mathrm D,*},\ldots,x_{N,\tau}^{\mathrm D,*}).
\]
The DFT barrier is
\[
    \Delta E_{\mathrm D}^{\ddagger}(\tau)
    =
    \max_{0\le i\le N}
    E_{\mathrm D}(x_{i,\tau}^{\mathrm D,*})
    -
    E_{\mathrm D}(x_{0,\tau}^{\mathrm D,*}).
\]

\paragraph{Fixed-path scoring.}
The fixed-path MLFF operator evaluates model energies on the same DFT-relaxed images:
\[
    \Delta E_{\mathrm{fixed}}^{\ddagger}(M,\tau)
    =
    \max_{0\le i\le N}
    E_M(x_{i,\tau}^{\mathrm D,*})
    -
    E_M(x_{0,\tau}^{\mathrm D,*}).
\]
It measures whether $M$ assigns the correct potential-energy profile to an externally supplied DFT path. The model does not choose or relax the path.

\paragraph{Self-relaxed MLFF NEB.}
The self-relaxed operator lets the MLFF define the forces that optimize the band:
\[
    \Gamma_{M,\tau}^{*}
    =
    \operatorname{NEB}_{M}(x_0,x_N)
    =
    (x_{0,\tau}^{M,*},\ldots,x_{N,\tau}^{M,*}),
\]
followed by
\[
    \Delta E_{\mathrm{self}}^{\ddagger}(M,\tau)
    =
    \max_{0\le i\le N}
    E_M(x_{i,\tau}^{M,*})
    -
    E_M(x_{0,\tau}^{M,*}).
\]
This operator measures optimization stability and qualitative path robustness on the MLFF's own PES. Because $\Gamma_{M,\tau}^{*}$ and $\Gamma_{\mathrm D,\tau}^{*}$ need not be the same curve in configuration space, $\Delta E_{\mathrm{fixed}}^{\ddagger}$ and $\Delta E_{\mathrm{self}}^{\ddagger}$ are not expected to agree. A low self-relaxed MLFF barrier is physically validated only after the path is checked against DFT, for example by DFT single-point energies on $\Gamma_{M,\tau}^{*}$ or by restarting a DFT NEB from the MLFF path.

\subsubsection{Error Decomposition}
\label{sec:error_decomposition_formal}

The fixed-path barrier error is
\[
    \epsilon_{\mathrm{fixed}}(M,\tau)
    =
    \Delta E_{\mathrm{fixed}}^{\ddagger}(M,\tau)
    -
    \Delta E_{\mathrm D}^{\ddagger}(\tau).
\]
Define the local energy residual and the model/reference maximum-image indices as
\[
\begin{gathered}
    \delta_M(x)=E_M(x)-E_{\mathrm D}(x),\\
    i_M^{*}=\arg\max_i E_M(x_{i,\tau}^{\mathrm D,*}),\qquad
    i_D^{*}=\arg\max_i E_{\mathrm D}(x_{i,\tau}^{\mathrm D,*}).
\end{gathered}
\]
The error can then be separated into endpoint and saddle-region terms:
\[
\begin{aligned}
    \epsilon_{\mathrm{fixed}}
    &=
    \underbrace{
    \delta_M(x_{i_D^{*}}^{\mathrm D,*})
    -
    \delta_M(x_{0}^{\mathrm D,*})
    }_{\text{endpoint-to-saddle energy calibration}}
    \\[2pt]
    &\quad+
    \underbrace{
    E_M(x_{i_M^{*}}^{\mathrm D,*})
    -
    E_M(x_{i_D^{*}}^{\mathrm D,*})
    }_{\text{maximum-image localization on the fixed path}}.
\end{aligned}
\]
Thus a model can fail by shifting the endpoint energy, miscalibrating the saddle energy, or moving the apparent maximum along the path. For a self-relaxed NEB, an additional geometric term appears because the model samples a different path:
\[
    \epsilon_{\mathrm{self}\rightarrow\mathrm D}(M,\tau)
    =
    \left[
    \max_{x_i\in\Gamma_{M,\tau}^{*}}E_{\mathrm D}(x_i)-E_{\mathrm D}(x_0)
    \right]
    -
    \Delta E_{\mathrm D}^{\ddagger}(\tau),
\]
which measures whether the MLFF-discovered path is also plausible on the DFT PES. In practice this term requires DFT single-point calculations or DFT NEB refinement on the MLFF-relaxed path.

\subsubsection{Robustness and Generalization Claims}
\label{sec:claims_formal}

The benchmark supports claims at three levels, with different evidence requirements.
\begin{enumerate}
    \item \textit{Pointwise accuracy}: $\mathcal E_{\mathrm{pt}}(M;\mu_{\mathrm{eq}})$ is small, as measured by energy/force RMSE, RDF agreement, and short-time stability.
    \item \textit{Process fidelity}: $|\Phi_M(\tau)-\Phi_{\mathrm D}(\tau)|$ is small for a specified task $\tau$, as measured by barriers, transport observables, sliding profiles, or DFT-path scoring.
    \item \textit{Robust generalization}: $\Phi_M$ remains accurate over a documented task family $\mathcal T$, with multi-path, multi-seed, OOD, and convergence checks.
\end{enumerate}
The present work claims Level 2 evidence for selected Cr-doped \ce{Sb2Te3} tasks and proposes Level 3 as a benchmark design principle requiring broader validation across materials, migration chemistries, architectures, and random seeds. This distinction prevents a successful result on one path from being over-interpreted as universal transferability.

\subsubsection{Benchmark Algorithm}
\label{sec:benchmark_algorithm_formal}

\begin{algorithm}[H]
\caption{Migration-centered MLFF benchmark}
\label{alg:migration_benchmark}
\begin{algorithmic}[1]
\State Define candidate models $\mathcal M=\{M_1,\ldots,M_L\}$ and reference method $D$.
\State Construct or audit equilibrium data $\mathcal S_{\mathrm{eq}}$ from DFT relaxations and AIMD.
\State Train or adapt models using documented training splits, seeds, and code hashes.
\State Measure pointwise errors $\mathcal E_{\mathrm{pt}}(M;\mathcal S_{\mathrm{eq}})$ and near-equilibrium observables.
\State Select migration and collective-displacement tasks $\mathcal T=\mathcal T_{\mathrm{ID}}\cup\mathcal T_{\mathrm{OOD}}\cup\mathcal T_{\mathrm{coll}}$.
\For{each task $\tau\in\mathcal T$}
    \State Generate endpoint-verified initial image sequence $\Gamma_\tau^0$.
    \State Compute or restart $\Gamma_{\mathrm D,\tau}^{*}=\operatorname{NEB}_{\mathrm D}(x_0,x_N)$ until convergence.
    \State Score each $M$ on the fixed DFT path to obtain $\Delta E_{\mathrm{fixed}}^{\ddagger}(M,\tau)$.
    \State Optionally run $\Gamma_{M,\tau}^{*}=\operatorname{NEB}_{M}(x_0,x_N)$ to assess self-relaxed stability.
    \State If $\Gamma_{M,\tau}^{*}$ differs from $\Gamma_{\mathrm D,\tau}^{*}$, validate the MLFF path with DFT single points or a DFT restart.
    \State Record barrier errors, endpoint errors, maximum-force histories, convergence flags, RMSD/SOAP overlap, and checksums.
\EndFor
\State Report model rankings with uncertainty over seeds, paths, and convergence status.
\State Interpret latent-space or SHAP diagnostics as supporting evidence, not as a replacement for task-level validation.
\end{algorithmic}
\end{algorithm}

\subsubsection{Applicability Assumptions and Acceptance Criteria}
\label{sec:acceptance_formal}

The framework is applicable when the process of interest can be represented by auditable endpoint states and a physically meaningful path or trajectory functional. It is especially appropriate for rare events, migration, sliding, stacking rearrangements, and other processes where a small set of high-leverage configurations controls the downstream observable. It is less complete when the dominant failure mode involves electronic excited states, bond-breaking chemistry absent from the reference functional, long-range electrostatics beyond the model architecture, or finite-size effects not represented in the chosen cell.

A benchmark task is accepted into the final manuscript only if the following checks are satisfied:
\begin{enumerate}
    \item DFT endpoints and NEB images are reconstructable from raw inputs and outputs.
    \item The DFT reference barrier is convergence-verified, or explicitly labeled as a current candidate value.
    \item Fixed-path and self-relaxed NEB quantities are reported as distinct operators.
    \item Model comparison includes seed sensitivity where training stochasticity can affect rankings.
    \item Training/test leakage is checked with grouped splits or path-aware proximity metrics.
    \item Path overlap is quantified with geometric RMSD and/or descriptor-space similarity.
    \item Any MLFF-discovered low-barrier path is tested on the DFT PES before being interpreted physically.
    \item Scripts, model identifiers, input files, convergence flags, and output checksums are preserved in the release manifest.
\end{enumerate}

\subsubsection{Physics and Data-Evidence Chain}
\label{sec:evidence_chain_formal}

The physical basis of the benchmark is transition-state theory: kinetic rates depend exponentially on activation free energies, so small barrier errors can lead to large rate errors,
\[
    k \propto \exp[-\Delta G^{\ddagger}/(k_{\mathrm B}T)].
\]
Consequently, a model that has low average force RMSE near equilibrium can still be unreliable if it misrepresents the PES along rare-event pathways. The evidence chain used here is:
\[
\begin{gathered}
\text{DFT/AIMD data}
\rightarrow
\text{trained MLFFs}
\rightarrow
\text{equilibrium validation}\\
\rightarrow
\text{path/process probes}
\rightarrow
\text{DFT ground-truth comparison}\\
\rightarrow
\text{representation diagnostics}
\rightarrow
\text{data-acquisition guidance}.
\end{gathered}
\]
This chain closes the benchmark loop: physical theory identifies the observables that matter, DFT supplies the reference PES, MLFFs are evaluated through measurement operators matched to downstream processes, and diagnostic analyses indicate which missing configurations should be acquired next.
